\section{Entities and Ideas}

In Lectures 5 and 6 of \emph{Divine Humanity}, \textbf{Vladimir Solovyov} explains the central points of his metaphysical understanding. Solovyov begins with Greek paganism, Neoplatonism, Hermetism, and German Idealism. However, Solovyov comes to a deeper conclusion and also adds some insights to \textbf{Rene Guenon}'s system, particularly in the \emph{Multiple States of Being}.

\paragraph{The Entity}

Although Solovyov claims that ``a continuous thread of logical thought must lead every consistent mind from the sensuous experience of phenomena to the speculative belief in ideas,'' we will simply stipulate that to be true as we have previously provided ample evidence. Once accepted, there are some obvious corollaries. The first is this:

\begin{quotex}
The usual thinking is that if all entities endowed with senses disappeared from the world, the world would nevertheless remain as it is, with all its variety of forms, with all its colors and sounds.
\end{quotex}


This is the naïve realist, and scientific, point of view. To understand its truth, all that is necessary is to try to imagine such a world. Of course it can't be done. Then he adds:

\begin{quotex}
The world we know is only a phenomenon for us and in us; it is our representation. To place it wholly outside ourselves, as something absolutely autonomous and independent of ourselves is a natural illusion.
\end{quotex}


Yet this representation is not arbitrary since we cannot create or destroy material objects at will. So the physical world has causes independent of us, and these are what Solovyov investigates. Solovyov calls these forces ``entities'', which have three qualities.

\begin{itemize}


\item \textbf{Atom}. An atom is an independent unit, the foundation of reality, real centers of being

\item \textbf{Monad}. A monad is a living, acting force. As such, monads interact with each other.

\item \textbf{Idea}. The idea is what distinguishes units from each other, i.e., the forces must represent a determinate content.
\end{itemize}


Although the translation uses the term ``the all'', I will instead call it the Infinite, or all-possibility, as more metaphysically precise in Englist. Solovyov understands this:

\begin{quotex}
The content of the Infinite consists of living and acting entities, eternal and abiding, which by their interaction form all reality, all that exists. ... The relation of each entity to the Infinite constitutes its \emph{objective} idea, which represents the full manifestation or actualization of its inner uniqueness, or \emph{subjective} idea.
\end{quotex}


These entities form a hierarchy, because any two entities themselves have a higher principle that is common to both of them. At the peak of the hierarchy is the idea of absolute goodness or absolute love. The fullness of the ideas cannot be understood as a ``heap'' or as their ``mechanical aggregate''. Rather their inner unity must be grasped, i.e., intuitively, and that is love.

Solovyov claims that philosophical systems are unbalanced when they focus on just one element---atom, monad, or idea---of the entity. What is required is the synthesis of all three.

\paragraph{Differentiation}


There are two opposite misunderstandings. The first is pluralism, or the notion that there are many independent entities. But how then could they interact without a common principle? The other error is to ignore the multiplicity of being and focus only on the higher elements of the hierarchy. The latter error is characteristic of the modern world, with its lack of differentiation. Solovyov writes:

\begin{quotex}
The more particular concepts a concept subsumes within itself, the fewer features it will have, the poorer its content will be, the more general, more indeterminate it will be. Thus, for example, ``humanity'' as a general concept embracing all human beings, and therefore of a scope much broader than the concept ``monk'', is much poorer than the latter in content.
\end{quotex}


Since such relations are the result of a process of abstraction, they have no content of their own. We see this today, for example, in those who claim to make no distinctions, but are rather ``philanthropists'', or ``lovers of humanity'', even if they love no one in particular. On the contrary, unlike abstractions, entities act upon each other, yet the focus on the abstraction hides the real, actual relations between the entities.

\paragraph{The Person as Entity}


Opposed to this abstraction, which is a negation of an entity's characteristic, is the idea as the positive determination of a particular entity. The entity then seeks to actualize its possibilities among other entities, not to strip away its defining character for the sake of some abstraction. Solovyov write:

\begin{quotex}
Every human person has a proper, individual character and represents a certain particular idea. By entering into interaction with others or in being determined by and determining others, every person discloses this proper, individual character, and realizes his own idea. Without this, this character and idea would remain pure possibility. They become actual only in the person's activity, in which the person is necessarily also determined by others. Consequently, this determination is not a negation, but an actualization.
\end{quotex}


\paragraph{Intellectual Intuition}


Next, Solovyov explains how we know these ideas. Since they are equally independent of both rationalist abstractions and sensuous reality, a special mode of thought activity is required, which is called intellectual intuition. This seems to be a stumbling block, since (1) there is the belief that we ``see'' things as they are or (2) that logical thinking is the way to grasp and manipulate ideas. However, neither option is possible since entities are living realities. These refutations may help:

\begin{enumerate}


\item Material reality represents in itself only contingent and transitory phenomena, but not autonomous entities or foundations of being.

\item Rational thinking is restricted to the abstractions, which have no positive content of their own. Hence, it cannot grasp the entities themselves.
\end{enumerate}


Sense experience and rational thought can only know accidental and particular facts, and not the necessary and universal truths or ideas. In short, he writes:

\begin{quotex}
Abstract thinking must serve either as an abbreviation of sense perception or as an anticipation of intellectual intuition, insofar as the general concepts forming it can be affirmed either as \emph{schemata of phenomena or as shadows of ideas}.
\end{quotex}


\paragraph{Postscript}


There is an exercise to make all this clear. Academic discussions of Solovyov will focus on his ``influences'', the relationship of his thought to other thinkers, a logical analysis of his ideas, etc., as though it were a construct to be analyzed. On the contrary, we use his writings as a series of guideposts for understanding and living in the world: to practice the development of intellectual intuition, to learn to ``see'' the entities, atoms, monads, as embodied ideas.


\flrightit{Posted on 2014-04-10 by Cologero}

\begin{center}* * *\end{center}

\begin{footnotesize}
\begin{sffamily}
\texttt{Mikkel on 2019-04-13 at 12:23 said:}

Reading this post and doing some of the practices, below are some interpretations of my experiences with it.

``Atom. An atom is an independent unit, the foundation of reality, real centers of being.'' -- Experience, consciousness, the direct channel of our science, the senses where we see what aligns and does not, the ``I'' we experience.

``Monad. A monad is a living, acting force. As such, monads interact with each other.'' -- Our story and lives, the interactions between others and the ways and whys influences bond and fade away between each other, A/B influences, law of affinity, the conscious real ``I''.

``Idea. The idea is what distinguishes units from each other, i.e., the forces must represent a determinate content.'' -- Our Being, Spirit and remembering our real Self, who and what we really are, the quality of our lives that is fixed but not restricting, the Self even beyond the Self that determines the responsibility of our creation.

\end{sffamily}
\end{footnotesize}

